\documentclass{article}

\usepackage[T1]{fontenc}

\usepackage{booktabs,graphicx,siunitx}

\begin{document}

\section{Summary statistics}

\begin{tabular}{lllllllllll}
\toprule
Group & Variable & N & M & SD & 95\% CI & Min & Max & Median & Skew & Kurtosis (excess) \\
\midrule
Female & exam\_score & 120 & 65.69 & 8.31 & [64.19, 67.20] & 44.84 & 84.12 & 66.27 & $-$0.17 & $-$0.55 \\
Male & exam\_score & 120 & 65.44 & 8.03 & [63.99, 66.89] & 44.91 & 81.23 & 65.83 & $-$0.19 & $-$0.56 \\
Female & study\_hours\_per\_week & 120 & 6.80 & 2.37 & [6.37, 7.22] & 1.84 & 11.44 & 6.67 & $-$0.07 & $-$1.11 \\
Male & study\_hours\_per\_week & 120 & 6.59 & 2.48 & [6.14, 7.04] & 2.14 & 11.62 & 6.45 & 0.11 & $-$1.17 \\
\bottomrule
\end{tabular}

\includegraphics{figures/01_summary-statistics/figure_histogram_exam_score.png}

\includegraphics{figures/01_summary-statistics/figure_histogram_study_hours_per_week.png}

Descriptives summarize your data: the mean and median show the center, the SD and range show the spread, and skew/kurtosis show how far the shape departs from a normal bell curve. There is no significance test here.

Table X reports descriptive statistics for the study variables.

Statistical basis: Descriptive statistics table (mean, SD, median, min/max, N, skew, kurtosis) (Revelle W (2024). psych: Procedures for Psychological, Psychometric, and Personality Research. Northwestern University, Evanston.).

\end{document}